\documentclass[11pt]{article}
\usepackage[utf8]{inputenc}
\usepackage[T1]{fontenc}
\usepackage{amsmath,amssymb}
\usepackage{booktabs}
\usepackage{geometry}
\geometry{margin=1in}
\usepackage{hyperref}
\hypersetup{colorlinks=true, linkcolor=blue, urlcolor=blue, citecolor=blue}
\usepackage{natbib}
\usepackage{xcolor}

\title{A Controlled Replication of a Self-Model Loop for Language-Model Agents:\\
Controls, Confounds, and Finite-Horizon Path Dependence\\
{\large Version 1.1 --- revised after external review}}

\author{Grant Wilson\thanks{Independent researcher. Correspondence via the
project repository.}}

\date{\today}

\begin{document}
\maketitle

\begin{abstract}
A recently described architecture wraps a large language model (LLM) with a
self-referential inference layer: a Bayesian belief over the agent's own
interaction ``stance'', re-injected into generation each turn, together with a
hidden deliberative monologue. The originating write-ups report strong effects
but do not publish the controls needed to separate the contribution of the state
\emph{content} from the presence of an instruction, an honest task baseline, or
the intrinsic inertia of the estimator. We give an independent, clean-room reimplementation of the loop's core (stance
tracker, state prefix, optional private monologue) and subject its central
claims to those controls across seven small experiments on one or a few model
endpoints.
We find: (i) under a structured-placebo control the state's \emph{content} is
behaviourally load-bearing only when the posterior is confident, and the private
monologue is the strongest single factor---though without a matched text-length
placebo we cannot isolate it as the unique causal channel; (ii) a clamped partial
factorial shows that mode-labelled prefixes reliably condition reply length and
question use (an instruction-following result, not a probabilistic calibration
test) and that the ``gate'' token is a substantial, non-uniform channel rather
than a minor one; (iii) on a ground-truth hidden-stance task the loop infers well
above chance (18/30; $p\approx2\times10^{-6}$) but shows no accuracy advantage
over an honest single-shot LLM baseline, which is itself far above chance; (iv)
driving the closed loop with a ramped numerical push produces substantial
finite-horizon path dependence with no observed sharp jump across seven sampled
levels; a state-decoupled null shows roughly two thirds of the raw separation is
accumulator arithmetic, leaving an arithmetic residual of $\sim$0.07 (mean
vertical separation) whose size and mechanism we do not claim to have
statistically identified; that residual did not collapse for dwell times up to
12 turns per level, but a 100-turn basin test shows relaxation still under way,
so slow lag is not excluded; and the basin test found trajectories from opposite
initial states converging and still closing at turn 100---no evidence of
bistability in this limited test, not a demonstration of a single attractor.
Every claim tested reproduces in a ``true-but-softer'' form once the missing
controls are added, and several of our own stronger intermediate readings failed
later controls. We make no claim about consciousness or any threshold measure;
the scope is strictly the loop's measurable behaviour. Code and data are
released. This is a preprint, not a peer-reviewed result.
\end{abstract}

\section{Introduction}

Adding a persistent, self-referential state to an LLM agent---a belief about the
agent's own stance, updated online and fed back into generation---is an
appealing idea with a growing informal literature. A prominent recent
instantiation (hereafter ``the AC1 loop'', following its author's naming)
combines four ingredients: a Bayesian stance tracker; a rendered state
\emph{prefix} injected upstream of the reply; a hidden \emph{private monologue}
that deliberates before the public reply; and a second-order term modelling the
interlocutor. The originating essays report large effects---for example that
removing the monologue changes the output on 92--100\% of measurable turns---and
frame these as evidence for strong properties.

Our interest is narrower and purely methodological. Effects of this kind can
arise for reasons that have nothing to do with a self-model: any prefix that
contains an imperative will change behaviour; a ``baseline'' set at chance will
make any above-chance system look impressive; and an estimator that accumulates
evidence will look path-dependent whether or not the surrounding loop is---with
decay, through inertia; without decay, by construction, since an undecayed
accumulator never forgets. The originating
write-ups do not report the controls that separate these. This paper supplies
them, on an independent reimplementation, and reports what survives.

\paragraph{Scope and non-claims.} We study only the measurable behaviour of the
loop. We make \textbf{no} claim about the $\Psi$-threshold, consciousness, or the
wider conjecture the architecture is embedded in; those questions are neither
tested nor implied here. Where a claim holds we say so; where it holds only in a
softer form, or rests on an unresolved confound, we say that too.

\section{The loop}

The state is a Dirichlet posterior over five stance modes
$\{$exploring, evaluating, overwhelmed, asserting, disengaging$\}$. Each turn:
\begin{enumerate}
\item a \textbf{sensor} (an LLM classification with confidence) reads the latest
line and yields a per-mode likelihood $\ell$ and a scalar evidence $e$;
\item the tracker updates $\alpha \leftarrow \rho\,\alpha + e\,\ell$, with a
mean-step cap, an $\alpha$ floor, and a concentration cap;
\item the posterior $p=\alpha/\kappa$ is rendered into a \textbf{prefix} (the
five values, the dominant mode, a \emph{gate} flag, and an optional per-mode
\emph{strategy} line);
\item an optional \textbf{private monologue} is generated from the prefix and
prepended, unseen, to the reply's system prompt;
\item the reply is produced; the sensor reads it; the loop repeats.
\end{enumerate}
The decay $\rho<1$ gives the posterior finite memory. Our reimplementation is
standard-library Python against any OpenAI-compatible endpoint; it copies no
part of the original codebase and is built only from the public descriptions.

\paragraph{A practical note that matters for reproduction.} Some reasoning
models return their text in a \texttt{reasoning\_content} field with an empty
\texttt{content}. A sensor reading only \texttt{content} then silently returns
nothing and the loop degrades \emph{without erroring}. Our extractor falls back
to \texttt{reasoning\_content}; anyone reproducing this against a reasoning model
must do the same or measure a broken loop.

\section{Experiments}

All behaviour metrics are ground-truth (word counts, question flags, matches to a
known hidden label); where a blind label is needed (E1) it is produced by a
separate model that never sees the prefix. Figures below are from a small number
of base models and should be read as illustrative in magnitude and indicative in
direction.

\subsection{E1: structured-placebo ablation}
Four arms at matched format: \textsc{true} (production prefix), \textsc{scrambled}
(same posterior with mode labels deranged, dominant and strategy rebuilt so the
placebo is internally consistent but wrong), \textsc{strategy\_stripped} (true
posterior with the strategy line removed; note it still carries the mode names
and values, the \texttt{Dominant}/\texttt{gate} fields and the ``internal
stance'' framing, so it is \emph{not} a pure numeric control---earlier versions
called it \textsc{numbers\_only}, which overstated what was stripped), and
\textsc{null} (no prefix). Under a \emph{weakly} held posterior,
\textsc{true} and \textsc{scrambled} were not separable on the blind-label
channel (JSD $0.007$ bits against a one-off split-half noise floor of $0.109$;
no formal equivalence test), although they differ visibly on length ($67.5$ vs
$46.3$ words), and \textsc{strategy\_stripped}${\approx}$\textsc{null}: the
imperative \emph{strategy line} carried the visible effect, not the state
values. Taken alone this is deflationary---the model follows the instruction
whether the state behind it is true or deranged. Scrambling also changes the
dominant label, gate text and strategy line together, so it does not isolate a
single content factor.

\paragraph{Reproducibility caveat.} The deposited E1 summary was produced by an
earlier variant of the script with a weakly held posterior; that variant's exact
posterior, generated replies and blind-judge outputs were not retained. The
released \texttt{e1\_placebo.py} clamps the dominant mode at $0.72$ and does not
reproduce the deposited E1 figures. E1 should therefore be read as illustrative;
the strong-state result rests on E2, which is reproducible from the released
code.

\subsection{E2: clamped factorial}
Clamping the posterior to each mode in turn (so the design cannot collapse into a
single mode) and toggling gate, monologue and strategy, we ran 30
configurations: 5 modes $\times$ gate $\times$ monologue with the strategy line
(20) plus 10 strategy-stripped configurations with the monologue off; each
configuration received 2 prompts $\times$ 3 repetitions, i.e.\ 180 replies in
total (an earlier version of this paper wrote ``180 cells''; that counted
replies, not configurations). The design is a \emph{partial} factorial: the
strategy-off\,$\times$\,monologue-on combination was not run, so the
strategy\,$\times$\,monologue interaction is not estimable. Note also that
``gate off'' inserts \texttt{gate=full scope} rather than omitting the field,
and the monologue condition appends the imperative ``Act on that assessment''
as well as the monologue text. Under a \emph{strong} clamped posterior:
\begin{itemize}
\item the dominant mode \textbf{predicts} reply behaviour (Table~\ref{tab:calib}).
This is measured with the strategy instruction \emph{on}, so it is largely
instruction following by construction; it is a mode-conditioned-style result,
not a probabilistic calibration test (no reliability curve, Brier or log-loss),
and we no longer call it ``calibration'';
\item \textsc{strategy\_stripped} prefixes now produce long, unshaped replies
($\sim$150--200 words) that differ sharply from strategy-bearing prefixes
(46--133 words): the state values are \emph{not} universally inert. This
\textbf{corrects} the E1 reading---inertness was a weak-state artefact;
\item the gate token moves length \emph{substantially} and non-uniformly:
disengaging $133.5\to56.0$ words ($-58\%$), asserting $46.0\to78.0$ ($+70\%$),
overwhelmed $108.0\to137.0$ ($+27\%$), six replies per mean. An earlier version
called this channel ``minor''; the data do not support that;
\item the monologue shortens the two longest states (disengaging
$133.5\to30.2$, overwhelmed $108.0\to76.0$) but \emph{lengthens} asserting
($46.0\to65.5$) despite the asserting strategy asking for one action and stop:
the effect is selective, not uniformly state-consistent.
\end{itemize}
The monologue is the strongest single factor in the factorial where the
clamped state is strongest (\texttt{e2\_factorial.json}). We stop short of
naming it the identified causal channel: the monologue condition also adds
text and an imperative, and we ran no matched placebo---equal-length irrelevant
text, a monologue generated from a scrambled state, the same monologue without
``Act on that assessment'', or monologues replayed across states. Without
those, the observed change could reflect extra instruction, semantic priming or
generic elaboration rather than processing of a valid self-state. The
originating write-ups claim the monologue is causally upstream; our data are
consistent with that but do not establish it. It is the leading candidate, not
a clean identification.

\begin{table}[t]
\centering
\begin{tabular}{lrr}
\toprule
Clamped mode & mean words & question rate \\
\midrule
exploring    & 51.0  & 1.00 \\
evaluating   & 102.2 & 0.50 \\
overwhelmed  & 108.0 & 0.83 \\
asserting    & 46.0  & 0.50 \\
disengaging  & 133.5 & 0.50 \\
\bottomrule
\end{tabular}
\caption{E2 mode-conditioned reply style: clamped dominant mode vs.\ reply
behaviour (strategy on, monologue off, gate off; six replies per row). The mode
predicts style---exploring is short and always questions; asserting is short
and direct; the deliberative modes run long. Because the strategy instruction
is on, much of this is instruction following.}
\label{tab:calib}
\end{table}

\subsection{E3: hidden-stance game with an honest baseline}
Each episode fixes a hidden mode; a simulator writes three short in-mode lines;
three predictors are scored against ground truth: random, an honest single-shot
LLM reading the whole transcript, and the sensor+filter loop. Over 30 episodes:
loop $0.60$ (18/30), single-shot $0.60$ (18/30), random $0.233$ (chance
$0.20$). Against fixed chance the loop is clearly above it (one-sided binomial
$p\approx1.8\times10^{-6}$), though the 95\% interval on $0.60$ is wide
(0.42--0.75). It shows no accuracy advantage over the honest single-shot
baseline; equal point estimates do not establish equivalence, and this design
does not measure persistence (that would need noisy, conflicting or time-varying
evidence), so ``the filter buys persistence rather than accuracy'' is a
hypothesis, not a finding of E3. The task is synthetic (an LLM told the mode
writes three lines; a classifier with closely related definitions labels them;
simulator, sensor and baseline default to the same base model), and the ``loop''
here is the per-line classifier plus accumulator, without the prefix-feedback
or monologue. The methodological point stands: a bare LLM is \emph{not} at
chance on this task, so reporting a chance-level baseline would inflate the
apparent gain of the loop.

\subsection{E4: closed-loop hysteresis, and the decay confound}
Only this experiment probes dynamics. Closing the loop and applying an external
push toward \emph{disengaging}, ramped up ($0\to0.9$) then back down
($0.9\to0$) from the same posterior, we measure the disengaging mass at each
level (Table~\ref{tab:hyst}). The down-ramp sits far above the up-ramp at every
level: the mean vertical separation across the seven sampled levels is $0.126$
(the result files call this \texttt{hysteresis\_loop\_area}; it is a fixed-grid
mean gap, not a geometric area---the trapezoidal area over push $0$--$0.9$ is
$0.116$). Pushed to disengagement at $0.9$ and released to $0$, the state
remains elevated over this horizon. The up-ramp rises smoothly (max single-step
jump $0.05$): no large jump is observed between the seven sampled levels, which
does not exclude a transition between them. Three design facts bound the
reading. The ``push'' is not an observation shown to the LLM: the code adds
$2.5\times$push to the disengaging component of the sensor likelihood, so the
estimator is driven numerically and history dependence in an accumulating
estimator is expected. Each curve is a single stochastic trajectory (sampling
temperature $0.7$) without repetitions or confidence intervals. And E4--E6
generate replies from the state prefix alone; they do not call the private
monologue, so the dynamics do not exercise the full architecture. So the loop
shows substantial \textbf{finite-horizon path dependence} but no observed sharp
transition.

\paragraph{The decay confound---and why the obvious control is insufficient.}
The tracker's decay ($\rho=0.97<1$) builds intrinsic inertia into the posterior,
so some stickiness is mechanical rather than emergent. The natural control is a
decay-free rerun: with $\rho=1.0$ the mean separation held at $0.115$ (vs.\ $0.126$)
---essentially unchanged (Table~\ref{tab:hyst}), and the transition remains
smooth under both settings. We initially read this as decisive. It is not: an
accumulator at $\rho=1.0$ never forgets, so it is path-dependent \emph{by
construction}, and ``survives decay-off'' cannot distinguish loop-borne memory
from accumulator arithmetic. Separating the two requires a null that keeps the
accumulator and removes the loop---which is E5b, below.

\begin{table}[t]
\centering
\begin{tabular}{lrrrrrrr}
\toprule
push & 0.00 & 0.15 & 0.30 & 0.45 & 0.60 & 0.75 & 0.90 \\
\midrule
up ($\rho{=}0.97$)   & .184 & .198 & .249 & .284 & .324 & .364 & .385 \\
down ($\rho{=}0.97$) & .382 & .405 & .417 & .421 & .419 & .422 & .401 \\
\midrule
up ($\rho{=}1.0$)    & .185 & .198 & .241 & .254 & .289 & .323 & .356 \\
down ($\rho{=}1.0$)  & .368 & .379 & .389 & .383 & .386 & .376 & .370 \\
\bottomrule
\end{tabular}
\caption{E4 disengaging mass along the up/down ramps (one trajectory each).
Mean vertical separation $0.126$ ($\rho{=}0.97$) vs.\ $0.115$ ($\rho{=}1.0$);
the separation is essentially unchanged with decay removed.}
\label{tab:hyst}
\end{table}

\subsection{E5--E6: what kind of memory?}
Three follow-ups decompose the E4 effect. All run at $\rho=1.0$ on a single
base model, with the push held for $K$ turns per level (E4 corresponds to a
fast ramp).

\paragraph{E5: ramp-rate sweep.} If the path dependence were lag---a viscous
state trailing a moving input---slowing the ramp far enough would collapse the
separation. Over $K=1,3,6,12$ turns per level it did not shrink
(Table~\ref{tab:memdecomp}). That is the supportable statement: \emph{no collapse
was observed for dwell times up to 12 turns per level}. It does not rule out
lag. E6 shows relaxation still under way after 100 turns, so $K\le12$ is far
from the quasi-static limit and the E5 pattern is compatible with very slow
convergence. Increasing $K$ also increases the total evidence delivered at each
level and, at larger $K$, reaches the estimator's concentration cap, so dwell
time and evidence dose are confounded in this sweep.

\paragraph{E5b: state-decoupled null.} The decisive control. The identical ramp
protocol is run with a likelihood that ignores the state entirely (no LLM in
the loop), so any loop area is pure accumulator arithmetic. The null still
produces areas of $0.10$--$0.15$ with the same growth in $K$
(Table~\ref{tab:memdecomp}): roughly two thirds of the raw E5 separation is
the estimator, confirming that E4's ``survives decay-off'' framing overstated
what that rerun showed. The closed-loop-minus-null residual is $+0.06$ to
$+0.08$ at each $K$ and did not shrink over this range. We report it as an
arithmetic residual, not a statistically identified causal effect: the null is
deterministic and not matched to the closed loop's mode and confidence
distribution, temporal correlation or API noise, and there is a single
closed-loop realisation per $K$ with no uncertainty estimate. A matched control
would preserve the real sensor outputs and break only the feedback
(replaying or permuting state-to-reply pairings) across many seeds.

\begin{table}[t]
\centering
\begin{tabular}{lrrrr}
\toprule
turns per level $K$ & 1 & 3 & 6 & 12 \\
\midrule
closed loop (E5)          & 0.161 & 0.210 & 0.213 & 0.231 \\
state-decoupled null (E5b) & 0.101 & 0.135 & 0.142 & 0.153 \\
\midrule
loop-attributable excess  & 0.060 & 0.075 & 0.071 & 0.078 \\
\bottomrule
\end{tabular}
\caption{Decomposing the path dependence ($\rho=1.0$; mean vertical
separation, one closed-loop trajectory per $K$). The null---the same accumulator
with the loop removed---carries most of the raw separation; the residual is
$\sim$0.07 at each $K$ and did not collapse over $K\le12$. It is an arithmetic
difference against an unmatched null, not a measured causal constant.}
\label{tab:memdecomp}
\end{table}

\paragraph{E6: two-point basin test.} Whether the residual memory amounts to
true bistability---two attractors at identical input---is tested by starting
the closed loop from opposite dominant modes under a clamped push of $0.45$.
At $N=30$ turns the trajectories still held a $\sim$0.20 gap, but the lower
trajectory was still climbing, so a 30-turn run cannot distinguish two
attractors from slow convergence to one (the shipped
\texttt{e6\_bistability.json} verdict, computed at $N=30$, should be
disregarded for exactly this reason). At $N=100$ the trajectories approach
each other ($\sim$0.60 vs $\sim$0.66 disengaging mass), with tail gaps of
$0.07$--$0.08$ and the lower trajectory still rising. This is consistent with
convergence to a common attractor but the run stops before convergence, and it
tests one push value, two initial states, one model and two trajectory pairs.
The supportable statement is \textbf{no evidence of bistability in this limited
test}, not that the system has a single attractor. It does revise E4's ``the
state stays in the basin'' reading: over a long horizon the state does not stay;
the return is slow.

\paragraph{Revised claim.} Under a cyclic numerical drive the closed loop
shows substantial finite-horizon path dependence, most of which is accumulator
arithmetic. A residual of $\sim$0.07 (mean separation) remains over an
unmatched null and did not collapse for dwell times up to 12 turns per level;
its size, mechanism and rate dependence are not established. Initial-condition
effects remain visible for over 100 turns with relaxation still under way, and
no bistability was observed in the one configuration tested. Every stronger
reading tested in this series, including two of our own along the way, failed
its control.

\section{Discussion}

Across these controls the picture is consistent: each tested claim reproduces in
a \emph{true-but-softer} form. In small synthetic experiments on one or a few
model endpoints, explicit mode-labelled prompt additions changed output length
and question use, most strongly when the posterior was confident and with the
private monologue as the leading---but not isolated---factor. A
classifier-plus-accumulator inferred synthetic modes well above chance without
an accuracy advantage over an honest single-shot baseline. Under a cyclic
numerical drive the closed loop produced substantial finite-horizon path
dependence, most of it accumulator arithmetic, with the size and mechanism of
any additional feedback contribution left uncertain; a 100-turn basin test found
no bistability and relaxation still in progress.
Two of our own intermediate readings (that a decay-free rerun was the decisive
confound control, and that a 30-turn basin gap indicated bistability) failed
their own follow-up controls and are corrected above; we report them because the
failure mode---an accumulator that is path-dependent by construction, and a
stability criterion that passes on slow drift---will face anyone reproducing
this loop.

None of this speaks to the stronger claims the architecture is often presented
with. We measured structural properties of a feedback loop; whether such
properties bear on any threshold or phenomenal question is outside what these
experiments can address, and we make no such claim. What we can say is narrower
and, we think, useful: the loop does several real things, and it does them in a
milder and better-controlled form than the original presentation asserts.

\paragraph{Limitations.} Figures come from a small number of base models; the
directions are the finding, not the magnitudes. The blind labeller in E1 has
limited power on the label channel (the word-count and question-rate channels
separate more cleanly). E3 uses short (three-line) episodes; longer evidence
might let the filter's persistence convert into an accuracy gain over the
single-shot baseline. The E5/E5b/E6 memory decomposition comes from a single
base model and one push geometry, from single stochastic trajectories without
repetition or confidence intervals; the $\sim$0.07 residual is an arithmetic
difference against an unmatched null and should not be read as a constant or as
an identified causal effect. E4--E6 do not exercise the private monologue. The
E2 factorial is partial and its controls are imperfectly named (see E1, E2).
The deposited E1 summary is not reproducible from the released E1 script. The
implementation is a conceptual partial reimplementation: it reconstructs the
stance tracker, state prefix and monologue from public descriptions but not the
second-order interlocutor term or the error-corrector component the originating
account describes, and the original code was unavailable for comparison, so
architectural fidelity is not verified. Matched placebos for the monologue, a
complete factorial, many seeds with uncertainty estimates, genuinely
probabilistic calibration tests, and dynamics run well beyond the observed
relaxation time are the natural next experiments.

\section*{Revision history}
\textbf{v1.1 (September 2026).} Revised after an external validity review of
v1.0. Corrections of fact: E2 is 30 configurations / 180 replies, not ``180
cells'', and is a partial factorial; the \textsc{numbers\_only} control is
renamed \textsc{strategy\_stripped} to reflect what it removes; the gate channel
is substantial (up to $\pm$60--70\% in length), not ``minor''; the E4/E5
``loop area'' is a mean vertical separation across sampled levels. Reproducibility:
the deposited E1 summary comes from an unreleased weak-clamp variant. Claims
softened to what the designs support: ``calibration'' becomes mode-conditioned
style; the monologue is the leading rather than identified causal channel;
E5's residual is an arithmetic difference that did not collapse for $K\le12$,
not rate-independent memory; E6 shows no evidence of bistability rather than a
single attractor; E3's ``persistence'' is a hypothesis. Title changed
accordingly. No data were re-run; all result files are those of v1.0 with
corrected notes.

\section*{Reproducibility}
A clean-room implementation of the loop and all experiments (E1--E6, including
the E5b state-decoupled null and the long-horizon E6 rerun), with the raw
result files, is released under the MIT license at the project repository. The
code uses only the Python standard library and targets any OpenAI-compatible
endpoint. The decay-free control is a single environment variable
(\texttt{SELFMODEL\_RHO=1.0}).

\section*{Acknowledgements}
The self-model loop architecture studied here is due to Lark Laflamme (RavenNest
Scientific). This is an independent reconstruction and set of controls; it
reproduces none of the original code and is built from the public descriptions.
We thank the author for write-ups clear enough to rebuild from.

\end{document}
